(1) 理想气体等温过程的做功 \(W_{h}=nRT_h\ln r\)、\(W_c=-nRT_c\ln r\)。两定容段 \(W=0\)，储热或放热量大小均是 \(Q_{\rm isoV}=nC_V(T_h-T_c)\)。若回热器有效率为 \(\eta_r\)，它在定容冷却时储存 \(\eta_r Q_{\rm isoV}\)，在加热时释放同样多的热量；其余 \((1-\eta_r)Q_{\rm isoV}\) 由外部热源与冷源承担。

(2) 对外部热端，除等温膨胀吸热 \(nRT_h\ln r\) 外，还要补足不完全回热的定容升温需求，所以 \(Q_h=nRT_h\ln r+(1-\eta_r)nC_V(T_h-T_c)\)。冷端排热 \(Q_c=nRT_c\ln r+(1-\eta_r)nC_V(T_h-T_c)\)。两式相减给净功 \(W=nR(T_h-T_c)\ln r\)，它不依赖理想化回热有效率，因为两定容段的内能变化互相抵消。

(3) 热效率 \(\eta=W/Q_h=\frac{nR(T_h-T_c)\ln r}{nRT_h\ln r+(1-\eta_r)nC_V(T_h-T_c)}\)。当 \(\eta_r=1\) 且热交换可逆时，效率退化为 \(1-T_c/T_h\)，即 Carnot 上界。回热不完全令分母增大而循环净功不变，效率严格降低。极限 \(r\to1^+\) 时净功趋零，不宜把形式上较差的效率混同于有用功率。

(4) 工质循环熵变严格为零，若外界只有 \(T_h,T_c\) 两热库，热库净熵增为 \(\Delta S_{\rm env}=-Q_h/T_h+Q_c/T_c=(1-\eta_r)nC_V(T_h-T_c)(1/T_c-1/T_h)\ge0\)。等温过程贡献互相抵消，回热损失造成正熵产生。式中假设热源冷源只接纳未回收的定容热；若另设有限温差传热过程，还会有其他额外不可逆熵产生。
